\documentclass[10pt,twocolumn]{article}
\usepackage[margin=0.75in]{geometry}
\usepackage{amsmath,amssymb,booktabs,graphicx,hyperref}
\title{LinUCB and Thompson Sampling for DeskDrop Article Recommendation, Evaluated via Replay and Inverse Propensity Scoring}
\author{Patrick Gomes de Oliveira}
\date{\today}
\begin{document}
\maketitle
\begin{abstract}
We recast CI\&T DeskDrop's logged article engagements as a contextual
bandit with a dynamically changing candidate arm set, evaluated with a
synthetic, explicit logging policy so that every round carries a known
propensity. Two classical linear bandits -- LinUCB~\cite{li2010linucb} and
linear Thompson Sampling -- share one coefficient vector across every
candidate article, with context built from a point-in-time user content
profile interacted multiplicatively with each candidate's TF-IDF vector.
We report both match-conditional replay evaluation and an unbiased
inverse-propensity-scored (IPS) off-policy value estimate, and show that a
naive concatenation of user and item features is insufficient for a linear
model to learn personalised relevance -- an explicit interaction term is
required.
\end{abstract}
\section{Introduction}
``crop-contextual-recommender'' studies a bandit with a small, fixed arm
set and agronomic hard constraints. This project studies the classical
news-recommendation setting LinUCB was originally introduced for (Li et
al., 2010): a \emph{dynamic} candidate pool that changes every round, which
rules out a per-arm posterior and requires a shared linear model scored
against per-candidate feature vectors instead. We additionally evaluate
off-policy via inverse propensity scoring, since DeskDrop -- like essentially
every real engagement log -- records only what a historical system showed
and whether the user engaged, never a counterfactual over what else could
have been shown.

\section{Dataset}
CI\&T DeskDrop (Kaggle ref
\texttt{gspmoreira/articles-sharing-reading-from-cit-deskdrop}) logs
person-article engagement events (view, like, comment, follow, bookmark)
with real timestamps, alongside an article catalogue (title, body text,
language, share time). Each interaction event becomes one bandit round:
the context is a person-specific, point-in-time profile and a candidate
article's content; the logged action is the article the person actually
engaged with.

\section{Feature engineering}
Features are declared in \texttt{conf/features.yaml}.
\texttt{article\_tfidf} is a 20-term TF-IDF vector fit on the full article
catalogue's title and body text -- catalogue-only, so it carries no
interaction-label signal. \texttt{user\_profile\_running\_mean} is the mean
content vector of a person's strictly-prior interactions, recomputed fresh
at every round with a strict ``$<$'' cut-off (point-in-time correct); a
person's first-ever round, or any round for a never-before-seen person,
receives the zero vector. \texttt{article\_popularity\_running} is the
strictly-prior interaction count for the round's own article, log-compressed,
and doubles as the simulated historical system's preference signal for the
logging policy (Section~4.3). A registry check asserts the YAML
specification and the implementation agree in both directions.

\section{Methodology}

\subsection{Context construction}
A naive context of $[\text{profile}_t,\ \text{article}_a]$ concatenated is
insufficient: a linear score $\theta^\top x$ over concatenated blocks has no
term that multiplies a profile coordinate by an article coordinate, so it
cannot represent "this article matches this person's topic interest" --
only independent, person-blind and article-blind additive effects. The
context actually used is

\begin{equation}
  x_{t,a} = \big[\,\text{profile}_t \odot v_a,\ v_a,\ \log(1+\text{pop}_{t,a}),\ 1\,\big],
  \label{eq:context}
\end{equation}

where $\odot$ is the elementwise product and $v_a$ is article $a$'s TF-IDF
vector. The interaction block $\text{profile}_t \odot v_a$ lets a single
shared linear coefficient per vocabulary term score exactly the weighted
topical overlap between the person and the candidate.

\subsection{LinUCB (shared parameters)}
A single ridge-regression posterior $A, b$ is shared across every
candidate: $A \mathrel{+}= x x^\top$, $b \mathrel{+}= r x$ after each
observed round, with UCB score $p(x) = x^\top \theta + \alpha\sqrt{x^\top
A^{-1} x}$, $\theta = A^{-1}b$. This is the disjoint-per-arm LinUCB of
\texttt{crop-contextual-recommender} collapsed to its shared-feature
special case, appropriate when arms are not fixed and do not individually
accumulate enough visits.

\subsection{Linear Thompson Sampling}
A conjugate Bayesian linear regression posterior~\cite{agrawal2013thompson}
$\mathcal{N}(\mu, \Sigma)$ over the same shared $\theta$ is updated as
$\Sigma^{-1} \mathrel{+}= xx^\top/\sigma^2$,
$\Sigma^{-1}\mu \mathrel{+}= rx/\sigma^2$. At each round a single sample
$\tilde\theta \sim \mathcal{N}(\mu, \Sigma)$ scores every candidate;
exploration falls out of posterior sampling rather than an explicit bonus.

\subsection{Synthetic logging policy and propensities}
DeskDrop logs no counterfactual: only the article actually shown (and
engaged with) is known. We assume the historical system surfaced articles
in proportion to their already-observed popularity, so at round $t$ a pool
of candidates is drawn (the logged article plus popularity-weighted
distractors) and each pool member's propensity is $\pi(a\mid x_t) =
\text{softmax}(\log(\text{pop}_{t,a}+\epsilon)/\tau)_a$. The logged
article's own propensity under this explicit policy is what IPS divides by.

\subsection{IPS off-policy evaluation}
For a frozen policy $\pi'$, the IPS estimator over held-out rounds is

\begin{equation}
  \hat V_{\text{IPS}}(\pi') = \frac{1}{N}\sum_{t=1}^{N}
  \frac{\mathbb{1}[\pi'(x_t) = a_t]\, r_t}{\pi(a_t \mid x_t)},
  \label{eq:ips}
\end{equation}

an unbiased estimate of $\pi'$'s expected reward under the logging
distribution when propensities are correctly specified, using every round
rather than discarding non-matching ones (replay evaluation~\cite{li2011unbiased}, by contrast,
estimates a conditional quantity and discards them).

\section{Experimental protocol}
Interaction rounds are sorted by timestamp and cut by quantile into
training (70\%, replayed online), validation (15\%) and test (15\%). We
compare LinUCB, linear Thompson Sampling, a popularity-greedy baseline, and
a context-blind random policy, reporting both match-conditional replay
evaluation (Section~4.5, discarding non-matching rounds) and the IPS
estimator (Equation~\ref{eq:ips}, using every round) on the test partition.

\section{Results}
\begin{table}[h]
\centering
\caption{Replay and IPS off-policy evaluation on the synthetic CI fixture (test partition).}
\label{tab:results}
\begin{tabular}{lrrrr}
\toprule
Model & Match rate & IPS estimate & ESS \\
\midrule
LinUCB            & 0.2952 & 3.0114 & 21.19 \\
ThompsonSampling  & 0.2190 & 2.1249 & 18.18 \\
PopularityGreedy  & 0.1048 & 0.6019 & 10.92 \\
Random            & 0.1143 & 1.0556 & 10.45 \\
\bottomrule
\end{tabular}
\end{table}

Table~\ref{tab:results} reports both evaluation protocols on the synthetic
CI fixture. LinUCB attains the highest replay match rate and the highest
IPS value estimate, with linear Thompson Sampling second; both clearly
beat the popularity-greedy and random floors. This confirms the context
design in Equation~\ref{eq:context}: an earlier version of this pipeline
that concatenated the raw profile and article vectors instead of their
elementwise product performed \emph{no better than random} regardless of
training length, because a linear model over concatenated blocks genuinely
cannot express personalised topical match. We report this as a concrete,
reproducible illustration of a common pitfall in linear contextual bandits
for recommendation, not merely as an implementation detail.

\section{Limitations and conclusion}
The logging policy and its propensities are synthetic: DeskDrop itself
records no propensity, so IPS here evaluates against an explicit, assumed
historical policy rather than the true unknown one, and the reported
effective sample size (Section~4.5) should be read as a diagnostic on
estimator variance, not a guarantee of correctness against reality. The
20-term TF-IDF vocabulary is deliberately small; a production system would
tune it against held-out log-likelihood rather than a fixed size. Posterior
mean serving (Section on serialization) discards the exploration component
of both bandits between retraining cycles, which is a standard and explicit
simplification for a batch-retrained system, not a limitation of the
underlying algorithms. Overall, both LinUCB and Thompson Sampling
substantially outperform non-contextual baselines once the context
correctly encodes a person-article interaction term.

\bibliographystyle{plain}
\bibliography{references}
\end{document}
